\chapter{Conclusións}
\label{chap:conclusions}

\lettrine{O} obxectivo principal deste traballo foi desenvolver unha solución móbil que facilitase o seguimento diario do tratamento anticoagulante con Sintrom, reducindo a dependencia da folla física e simplificando tanto a interpretación da pauta como o control das tomas. Para alcanzar este obxectivo desenvolveuse SintromApp, formada por un cliente móbil implementado con Flutter e unha \gls{api} desenvolvida con FastAPI, complementados con diferentes servizos para o procesamento documental, as notificacións e a comunicación entre dispositivos.

Unha das partes máis relevantes do traballo foi o estudo das alternativas para automatizar a lectura das follas de tratamento. Ao longo do proxecto avaliáronse diferentes técnicas de recoñecemento óptico de caracteres, visión artificial e procesamento documental. As probas realizadas mostraron que non era suficiente con recuperar correctamente o texto presente no documento, xa que resultaba imprescindible conservar tamén a relación entre cada data e a súa dose para poder reconstruír a pauta de forma fiable. Finalmente, a comparación realizada permitiu seleccionar o modelo \textit{Template} de Azure AI Document Intelligence, que mostrou un comportamento máis estable na extracción da estrutura das follas e, especialmente, da táboa de doses.

A partir desta extracción implementouse o procesamento necesario para transformar a información obtida nunha pauta estruturada que puidese ser empregada directamente pola aplicación. Para iso foi necesario interpretar datas e doses, identificar os días de control, resolver cambios de ano, comprobar determinados valores e ordenar cronoloxicamente o calendario. Este proceso permitiu pasar dunha folla pensada para a súa consulta visual a unha representación dixital apta para o seguimento diario do tratamento.

En relación cos obxectivos formulados ao inicio do proxecto, considérase que estes se cumpriron de maneira satisfactoria. A aplicación permite incorporar unha folla mediante unha fotografía, unha imaxe almacenada no dispositivo ou un documento \gls{pdf}, reducindo a necesidade de introducir manualmente a información. Os datos extraídos poden revisarse e corrixirse antes de seren incorporados como pauta activa, o que resulta especialmente importante ao tratarse de información relacionada coa medicación. Unha vez gardada, a persoa usuaria pode consultar a dose correspondente a cada día, acceder ao calendario e ao histórico e visualizar a evolución do tratamento.

Tamén se incorporaron mecanismos destinados a facilitar o cumprimento da pauta. A aplicación permite configurar unha hora habitual para a toma e programa recordatorios para as doses pendentes. A persoa usuaria pode confirmar cada toma e o sistema conserva o seu estado, distinguindo entre tomas realizadas, realizadas fóra da hora prevista, pendentes ou non realizadas. Ademais, incorporouse un modo sinxelo co obxectivo de reducir a complexidade da interface nos casos nos que se precise unha interacción máis directa coa aplicación.

Outro dos obxectivos consistía en facilitar a supervisión do tratamento. Para iso desenvolveuse un sistema de vinculación entre pacientes e persoas supervisoras mediante códigos \gls{qr}. Esta relación permite compartir cambios relevantes da pauta e do estado das tomas entre dispositivos e posibilita tamén que unha persoa supervisora incorpore unha nova folla en nome dun paciente vinculado. Deste xeito, o seguimento non queda limitado ao propio dispositivo do paciente e pode realizarse tamén de forma remota.

A privacidade da información foi outro dos aspectos considerados durante o deseño e a implementación. Os datos do tratamento almacénanse principalmente de forma local e a \gls{api} non mantén unha base de datos central coa información clínica das persoas usuarias. A configuración sensible e as claves asociadas ás vinculacións almacénanse mediante mecanismos de almacenamento seguro, mentres que a información intercambiada entre dispositivos se cifra antes do seu envío. Desta forma, os servizos empregados para transportar as mensaxes non necesitan interpretar o contido clínico transmitido.

A solución desenvolvida foi comprobada mediante probas unitarias, de integración e de interface, complementadas con probas manuais en dispositivos Android. A \gls{api} alcanzou unha cobertura global do 96,1\,\%, mentres que o cliente móbil obtivo un 72,8\,\%. A diferenza entre ambos valores débese principalmente ás funcionalidades do cliente que dependen directamente do dispositivo, como a cámara, as notificacións locais ou determinadas operacións de sincronización, que foron verificadas tamén mediante probas manuais. Estas comprobacións permitiron validar os principais fluxos da aplicación, incluíndo a incorporación e revisión dunha folla, a persistencia da pauta, a confirmación das tomas, a vinculación entre dispositivos e a recepción das notificacións.

Durante o desenvolvemento foi necesario afrontar diferentes dificultades técnicas. A principal estivo relacionada coa variabilidade das condicións nas que pode capturarse unha folla e coa necesidade de reconstruír correctamente a información dunha táboa a partir dos datos extraídos automaticamente. A avaliación das distintas alternativas mostrou que unha boa taxa de recoñecemento de caracteres non implica necesariamente que unha ferramenta sexa adecuada para este problema, xa que a conservación da estrutura do documento resulta fundamental. Outro reto importante foi coordinar a persistencia local, os recordatorios e a sincronización entre dispositivos sen depender dun servidor central que almacenase a información do tratamento.

A solución presenta tamén limitacións que deben terse en conta. O sistema de extracción foi deseñado e probado sobre o formato das follas de tratamento empregadas polo SERGAS, polo que non pode asumirse o mesmo comportamento con documentos doutros organismos ou cunha estrutura diferente. A calidade da captura pode afectar igualmente ao resultado da extracción e, por este motivo, a aplicación incorpora unha etapa de revisión antes de gardar a información. Ademais, aínda que se realizaron probas funcionais e comprobacións de interface, non se levou a cabo un estudo formal de usabilidade con participantes externos nin unha validación clínica por parte de profesionais sanitarios. Estas comprobacións serían necesarias antes de considerar a utilización da aplicación nun contexto real.

Como posibles liñas de traballo futuro, poderíase ampliar a extracción para admitir outros formatos de follas de tratamento, incrementar a robustez ante documentos capturados en condicións desfavorables e continuar mellorando os mecanismos de sincronización entre dispositivos. Tamén resultaría de interese realizar estudos de usabilidade con pacientes e persoas supervisoras, especialmente con persoas de idade avanzada, co obxectivo de avaliar o modo sinxelo e detectar posibles barreiras de interacción. Finalmente, unha evolución máis ambiciosa da solución podería estudar a integración cos sistemas oficiais de saúde.

En conxunto, o proxecto permitiu desenvolver unha solución funcional que integra desenvolvemento móbil, deseño e implementación de \glspl{api}, procesamento documental, persistencia local, cifrado, mensaxería e probas. Máis alá do resultado técnico, o traballo permitiu comprobar a importancia de avaliar experimentalmente as tecnoloxías antes de incorporalas á solución definitiva e de adaptar as decisións de deseño ás características concretas do problema. SintromApp constitúe así unha base funcional sobre a que se podería continuar traballando para avanzar cara a unha ferramenta máis completa de apoio ao seguimento do tratamento anticoagulante.